% 本文件由 LaTeX 排版 Agent 根据有效 translation.json 自动生成。
% author=Tertullian; work_id=0409; title=论逃避迫害

\chapter{论逃避迫害}

% structure: p0001 source=h1 target=omitted reason=page-title-already-rendered-by-parent

1. 我的兄弟\Person{法比乌斯}，你最近曾问，因一些消息或别的什么被传开，我们在迫害中是否应当逃避。至于我，当时当场作了一些与地点和时间相宜的否定性观察，也因某些人的粗鲁，只带着未完全论述的题目离开，打算现在用我的笔更充分地论述它；因为你的询问使我对此感兴趣，而时代的处境本身也早已催逼我思考。随着越来越多的迫害威胁我们，我们就越发被召唤去认真思考这个问题：信德应当如何接受它们，而仔细思考的责任同样关乎你，你无疑因不接受护慰者——一切真理的引导者——自然地至今也曾在其他问题上反对我们。因此，我们也对你的询问作了有条理的处理，因为我们看到我们必须首先就迫害本身的情况作出决定：它是来自天主还是来自魔鬼，这样我们才能较不困难地站稳关于应对它的责任；因为凡事当知道了它的来源时，我们的认识就更清楚。的确，只断言（排除一切其它）没有什么事不发生在天主的旨意之下就足够了。但是，为了避免我们偏离当前问题，我们不会因这一断言而立刻给其他讨论提供机会，若有人回答说：\Quote{因此，邪恶和罪都来自天主；魔鬼从此，甚至我们自己，都完全自由。}当前的问题是迫害。关于此，我暂且说：没有什么事不发生在上天的旨意之下；理由是迫害特别相配于天主，并且可说是必需的，用于考验，或者如果你愿意，用于拒绝祂的宣认仆人。因为迫害的结果是什么，除了信德的考验和拒绝之外，还有什么其他结果？在这方面，主当然会筛分祂的子民。迫害，借着它一个人被宣布为合格或被拒绝，正是主的审判。但审判理当唯独属于天主。这就是那簸箕，它甚至现在就在清洁主的打谷场——我指的是教会——簸扬那混杂的信徒堆，从否认者的糠秕中分离出殉道者的麦粒；这也是\Person{雅各伯}所梦见的天梯，上面看到一些人升到更高处，另一些人降到更低处。同样，迫害也可被视为一场竞赛。这场比赛由谁宣告？唯有那位提供冠冕和奖赏者。你在\BibleBook{}{默示录}中找到它的谕令，列出祂激励人得胜的奖赏——尤其是那些在迫害中争战得胜的人，他们实际上不是与血肉争战，而是与邪恶的属灵势力争战。同样，你也会看到，裁判比赛也属于同一位荣耀的主，作为裁判者，祂召唤我们领取奖赏。迫害中的一件大事是促进天主的荣耀，正如祂试验和抛弃，施加和解除。但是，凡关乎天主荣耀的事，必然凭祂的旨意成就。对天主的信赖何时比在更大的敬畏中，以及当迫害爆发时更强烈？教会肃然敬畏。那时，信德在预备上更热切，在禁食、聚会、祈祷、谦卑、兄弟友爱和爱德、圣洁和节制上更有纪律。实际上，除了敬畏和盼望，别无他物。所以，正因如此，我们清楚证明了迫害，因为它改善天主的仆人，不能归咎于魔鬼。\par

2. 若不义并非来自天主，而是来自魔鬼，且迫害包含不义（因为还有什么比对待真天主的众主教、一切真理的追随者如同最卑劣的罪犯更不义的呢？），如此看来，迫害似乎来自魔鬼，因构成迫害的不义由它施行。我们应当知道——正如你们既没有不藉魔鬼之不义的迫害，也没有不藉迫害之信德的考验——为考验信德所必需的不义，并非使迫害合法化，而是提供了执行工具；实际上，就作为迫害理由的信德考验而言，天主的旨意先行，但就作为考验途径的迫害工具而言，魔鬼的不义随后。因为在其他方面，不义愈是对义德显出其敌意，便愈为它所敌对者提供证明的机会，使义德在不义中得以成全，正如力量在软弱中得以成全。原来天主拣选了世上软弱的，以羞辱那强壮的；拣选了世上愚拙的，以羞辱那聪明的。如此，不义也被利用，使义德因羞辱不义而得到认可。因此，既然这服务并非出于自由意志，而是出于服从（因为迫害是主为考验信德所规定的，但其执行则是魔鬼的不义，为使迫害得以兴起），我们相信迫害确实藉魔鬼之手发生，但并非源自魔鬼的主动。撒殚不得自由地反对永生天主的仆人，除非主准许，或是为了使祂藉着在考验中得胜的选民的信德击败撒殚自身，或是在世人面前表明那些背教投靠魔鬼的人实际上一直是祂的仆人。你们有\Person{约伯}的例子：除非从天主那里获得权柄，魔鬼不能以考验加害于他，甚至不能触及他的财产，除非主曾说：\BibleQuote{看，凡他所有的，我都交在你手中；只是不要伸手加害他本身。}简言之，若非后来因它的请求，主也准许了这一点，说：\BibleQuote{看，我把他交给你；但要保存他的性命。}同样，魔鬼也寻求机会试探宗徒们，只能凭特准，因主在\WorkTitle{福音}中对\Person{伯多禄}说：\BibleQuote{看，撒殚求得了许可，要筛你们像筛麦子一样；但我已为你祈求，叫你的信德不至丧失。}这就是说，魔鬼没有得到足以危及他信德的权柄。由此显明，动摇信德与保护信德二者都属于天主，因为二者都从祂那里求得——动摇由魔鬼求，保护由圣子求。当然，既然信德的保护已完全交托给天主子，祂从父那里祈求（父将天上地上的一切权柄都赐给了祂），那么魔鬼如何能凭\Emph{他}自己的权柄攻击信德呢？但在我们所领受的祈祷文中，当我们对父说：\BibleQuote{不要让我们陷于诱惑}（还有什么诱惑比迫害更大呢？），我们承认这临到我们的事，是按祂的旨意，我们祈求祂免除我们。因为这句话之后紧接着：\BibleQuote{但救我们免于凶恶}，意即不要将我们交给那恶者而使我们陷于诱惑，因为当我们没有被交给他受试探时，我们便从魔鬼的权柄下被解救出来。若非从天主那里获得准许，魔鬼的军团连对猪群也无权；何况对天主的羊群呢？我甚至可以说，当时猪的鬃毛也被天主数过，更不用说圣人的头发了。必须承认，魔鬼对于那些不属于天主的人，似乎确实有权柄——在此意义上确实是它自己的——因为外邦人被天主一次算作桶里的一滴水、禾场上的灰尘、口中的唾沫，从而被抛给魔鬼，如同自由产业一般。但那些属于天主家里的人，它不能凭任何自己的权利触动，因为圣经所载的案例表明何时——即为何原因——它可以触动他们。或者，为使众人得以被认可，试探的权柄被赐给它，或是应它的挑战，或是它被挑战，如已提及的诸例；或者，为相反的结果，罪人被交付给它，如同交给属于惩罚之执行的刽子手，如\Person{撒乌耳}的例子。经上记着：\BibleQuote{上主的神离开了\Person{撒乌耳}，有恶神从上主那里来扰乱他、窒息他。}或者目的是为了使人谦卑，如宗徒所说：\BibleQuote{有一根刺加在我身上，就是撒殚的使者来攻击我。}就连这一类事，若不同时使忍耐的力量在软弱中得以成全，也不得加于圣徒身上。因为宗徒也曾把\Person{非革罗}和\Person{赫摩革乃}交付给撒殚，使他们受管教而学会不再褻瀆。你们可见，魔鬼甚至从天主仆人那里获得权柄更为恰当；它凭自己的权利拥有权柄是绝不可能的。\par

3. 因此，既然看到这些情况在迫害中比其他时候更常发生，因为那时我们中间有更多的考验或拒绝，更多的辱骂或惩罚，那么必然地，它们的普遍发生是那位允许或命令它们即使是局部发生者所允许或命令的；我是指那一位，祂说：\BibleQuote{我是造和平又造邪恶的}——也就是战争，因为那是和平的对立面。但我们的和平还有什么别的战争，除了迫害？如果在其结果上，迫害强有力地带来生或死，创伤或医治，那么你也拥有这位作者。\BibleQuote{我要击打也要医治，我要使人活也要使人死。} \BibleQuote{我要焚烧他们，}祂说，\BibleQuote{如同金子被焚烧；我要试炼他们，}祂说，\BibleQuote{如同银子被试炼，}因为当迫害的火焰正在消耗我们时，我们信德的坚毅便被证明。这些将是魔鬼的火箭，信德借此得到焚烧和点燃的职务；然而这是出于天主的旨意。关于这一点，我不知道谁会怀疑，除非是那些信德轻浮而冷漠的人，这种信德抓住那些战战兢兢聚集在教会里的人。因为你说，既然我们毫无秩序地聚集，同时聚集，成群地涌入教会，外邦人便会被引导来查问我们，我们害怕会激起他们的焦虑。难道你不知道天主是万有的主吗？如果是天主的旨意，那么你将遭受迫害；如果不是，外邦人就会安静。要最确定地相信，如果你确实相信那位天主，没有祂的旨意，连一文钱买得的麻雀也不会掉在地上。但我想，我们比许多麻雀更贵重。\par

4. 那么，若从谁那里发出迫害是明显的，我们就能立即满足你们的疑惑，并单从这些引论就断定：人不应在其中逃跑。因为若迫害出自天主，我们绝无义务逃离以天主为作者的事；有双重理由反对：因为出自天主的事，一方面不应避免，另一方面也无法逃避。不应避免，因为它是善的；因为凡天主所注视的一切，必定是善的。也许\BibleBook{}{创世纪}中的这句话正是出于此意：\BibleQuote{天主看了，认为好}；并非祂若不看就不知道其善，而是以此表述来表明，它之所以是善，是因为被天主所观看。确实有许多事件是出于天主的旨意而发生，并且对某人有害。尽管如此，一件事物之所以为善，乃因它来自天主，是神圣的、合理的；因为什么是神圣的，而又是非合理、非善的呢？什么是善的，而又非神圣的呢？但如果对于人类的普遍认知而言，情况看似如此，那么在判断时，人的认知能力并不预决事物的本性，而是事物的本性预决其认知能力。因为每一个别的本性都是一个确定的实在，它要求知觉能力按其本来的样子去知觉它。现在，若来自天主的在其自然状态中确实是善的（因为凡来自天主的没有不善的，因为它是神圣的、合理的），而仅仅在人的能力看来是恶的，那么对于前者一切都是正确的；过错在后者。贞洁在其真实本性上是极善的，真理和正义也是如此；然而许多人却厌恶它们。也许真实本性因此而被牺牲给知觉感？这样，迫害在其自身本性上也是善的，因为它是一个神圣而合理的安排；但那些遭受它作为惩罚的人并不觉得它愉快。你看见，既然是出于祂，即使那恶也有一个合理的根据，当一个人在迫害中被逐出救恩状态时，正如你看见对善也有一个合理的根据，当一个人因迫害而使其救恩变得更加稳固时。除非，如同取决于主，一个人要么不合理地丧亡，要么不合理地被救，他将不能把迫害说成是恶的，因为这迫害在理性的指导下，即使在其恶的方面也是善的。所以，若迫害无论如何都是善的，因为它有一个自然的基础，我们就有充分理由断定：善不应被我们回避，因为拒绝善是罪；此外，凡已被天主所注视的，实际上已不能再被避免，既是出于天主，从祂的旨意逃脱是不可能的。因此那些认为自己应该逃跑的人，要么指责天主行恶，如果他们逃避迫害如避恶（因为无人回避善）；要么自认为比天主更强：他们这样想，以为当天主乐意这些事发生时，他们有可能逃脱。\par

5. 但有人会说：\Quote{我逃跑是我该做的事，以免我因背弃而丧亡；至于祂，若祂愿意，可以在我逃跑时带我回到审判台前。}首先回答我：你确定若不逃跑你就会背弃，还是不确定？因为如果你确定，你已经背弃了，因为你预设你会背弃，你就已经将自己交给了你所预设的那个结果；现在你想借着逃跑来避免背弃，那是徒然的，因为你心中已经背弃了。但如果你对此怀疑，为什么在恐惧的不确定中，在两种不同结果间摇摆，你不假定自己更能扮演宣信者的角色，从而增加你的安全，使你无需逃跑，正如你预设背弃会使你成为一个逃跑者一样？事情是这样：要么这两件事都在我们自己的能力内，要么它们完全由天主决定。如果宣认或背弃在于我们，为何我们不选择更高贵的事，即我们要宣认？如果你不愿意宣认，你就不愿意受苦；而不愿意宣认就是背弃。但如果事情完全在天主手中，为何我们不将之交托于祂的旨意，承认祂的权能与力量，正如祂能在我们逃跑时将我们带回审判，祂也能在我们不逃跑时保护我们；是的，甚至当我们生活在人群之中？这是奇怪的行为，不是吗？在逃避迫害的事上尊敬天主，因为祂能将你从逃跑中带回审判台前；但在作证的事上，却因怀疑祂的手能保护你免受危险而大大羞辱祂？你为何不在这边，即坚持与信靠天主的一面，说：\Quote{我尽我的本分；我不离开；天主若愿意，祂自会保护我。}？我们最好是留在原位，顺从天主的旨意，而不是按自己的意愿逃跑。\Person{鲁提略}，一位圣洁的殉道者，屡次从一个地方逃到另一个地方以逃避迫害，甚至用金钱为自己买来他认为的安全，然而尽管他为自己准备了彻底的安全，最终还是出乎意料地被捕了；他被带到长官面前，遭受酷刑，被残忍地折磨——我相信这是对他逃跑的惩罚——之后被投入火焰，如此将他曾回避的苦难偿还给了天主的仁慈。主藉此例子还想告诉我们什么，岂不就是我们不应逃避迫害，因为若天主不认可，逃避毫无益处？\par

6. 不，有人说，他实行了命令，当他从一城逃到另一城时。因为如此，有一个人——同样也是个逃亡者——选择这样持守，其他人也照样做了，他们不愿理解主那番宣告的意义，好把它用作他们怯懦的遮盖，尽管它有它特别适用的人、时间和理由。祂说：\BibleQuote{当他们开始迫害你们时，要逃往别的城。}我们坚持认为，这特别属于宗徒本人以及他们的时代和环境，正如随后的句子将显示的，那些句子只适用于宗徒：\BibleQuote{不要走外邦人的路，不要进入撒玛黎雅人的城，宁可往以色列家迷失的羊那里去。}但对我们来说，外邦人的路也敞开了，因为我们实际上正是在外邦人中找到的，并且一直走到最后；没有一座城被除开。所以我们向全世界宣讲；不但如此，连对以色列的特殊责任也没有交托给我们，除非我们也必须向万民宣讲。是的，如果我们被抓，我们不会被带到犹太人的公议会，也不会在犹太人的会堂里被鞭打，而是肯定会被带到罗马官长和审判台前。所以，宗徒的环境甚至要求逃跑的命令，他们的使命是首先向以色列家迷失的羊宣讲。因此，为了使这个宣讲在那些首先应该接受的人中得以完全实现——好让儿子们在狗之前得到饼，为此祂命令他们那时暂时逃跑——不是以躲避危险为目的，严格来说，在迫害催逼的借口下（相反，祂惯于宣告他们将受迫害，并教导这些必须忍受）；而是为了促进福音信息的宣扬，免得他们一旦被压制，福音的传播也受到阻碍。他们也不是要偷偷地逃到哪座城，而是像到处要宣讲他们的信息一样；并为了这目的，到处要受迫害，直到他们完成他们的教导。因此，救主说：\BibleQuote{你们还没有走遍以色列的城。}所以逃跑的命令局限于犹太地区。但没有指出犹太是特别宣讲范围的命令适用于我们，因为现在圣神已倾注在一切血肉之人身上。因此，\Person{保禄}和宗徒们自己，记念主的诫命，在以色列面前作这庄严的见证——以色列现在已被他们的道理所充满——说：\BibleQuote{天主的道本来必须首先传给你们；但你们既然弃绝它，并断定自己不配得永生，看，我们转向外邦人。}从那时起，他们就转身离去，正如他们前面的人所规定的，前往外邦人的路，进入撒玛黎雅人的城；以至于事实上，他们的声音传遍全地，他们的言语传到地极。因此，如果禁止踏入外邦人的路、进入撒玛黎雅人城的禁令已经终止，那么同时发出的逃跑的命令为什么没有终止呢？所以，从以色列满了之后，宗徒转向外邦人之时起，他们就不再从一城逃到另一城，也不迟疑去受苦。不，连\Person{保禄}——他曾因从城墙上被缒下而逃脱迫害，因为这样做在当时是命令的事——现在在他职事的末期，在命令已经终止之后，同样拒绝顺从门徒们的焦虑，他们恳切地求他不要冒险去\Place{耶路撒冷}，因为\Person{阿加波}预言了等待他的苦难；但\Person{保禄}做了完全相反的事，他这样说：\BibleQuote{你们为什么这般哭泣，使我心碎呢？我为主耶稣基督的名不但愿意被捆绑，也愿意死在耶路撒冷。}于是他们都说了：\BibleQuote{愿主的旨意成就。}主的旨意是什么？当然不再是逃避迫害。否则那些希望他逃避迫害的人，也可以引证主先前的旨意，其中他命令逃跑。所以，即使在宗徒们自己的日子里，逃跑的命令也是暂时的，就像当时同时命令的其他事一样，那命令不能继续与我们同在，因为它与我们的导师们一同终止了，即使它当初不是特别为他们发出的；或者如果主希望它继续，那宗徒们就做错了，因为他们没有谨慎地一直逃到最后。\par

7. 现在让我们看看，主其余的法令是否也与持久逃跑的命令一致。首先，的确，如果迫害来自天主，我们被命令避开它，而命令者正是那将它降在我们身上的那一位，这该作何理解？因为如果祂想要它被避开，祂最好不把它送来，免得祂的旨意看上去被另一个旨意挫败。\par

因为祂愿意我们或遭受迫害，或逃避迫害。若应逃避，如何遭受？若应遭受，如何逃避？事实上，命令人逃避又敦促人遭受——这完全相反——的旨意何其矛盾！\BibleQuote{凡在人前宣认我的，我在我天上的父前也必宣认他。} 他若逃避，如何宣认？他若宣认，如何逃避？\BibleQuote{凡以我为羞耻的，我在我天上的父前也必以他为羞耻。} 我若规避受苦，就是羞于宣认。\BibleQuote{为我的名字遭受迫害的人是有福的。} 因此，那些因逃跑而不按神圣命令受苦的人是有祸的。\BibleQuote{那坚持到底的，必要得救。} 那么，当你命令我逃跑时，如何又愿你坚持到底？若如此彼此对立的观点不符合天主的尊荣，它们便清楚证明：逃跑的命令在发出时自有其理由，我们已指出。但有人说，主顾念祂子民中一些人的软弱，出于慈爱，也向他们提供了逃跑这一避难所。难道祂竟不能——不用逃跑这等卑劣、不相称且奴性的庇护——在迫害中保守那些祂知道是软弱的人吗？事实上，祂并不珍视软弱者，反倒常常弃绝他们，祂首先教导的，不是我们要逃避迫害者，而是不要惧怕他们。\BibleQuote{你们不要害怕那能杀害肉身，却不能杀害灵魂的；但更要害怕那能使灵魂和肉身都陷入地狱的。} 然后祂为胆怯者分配了什么？\BibleQuote{那爱惜自己生命胜于我的，不配作我的门徒；不背起自己的十字架跟随我的，也不能作我的门徒。} 最后，在 \BibleBook{}{默示录} 中，祂没有为\BibleQuote{胆怯的}提供逃跑，而是在火与硫磺的湖——即第二次死亡——中，与其余被弃者一同承受悲惨的分。\par

8. 祂有时也亲自逃避暴力，但原因与祂命令宗徒们逃避的相同：即祂要完成教导的职事；当这职事完成时，我不说祂站立稳固，但祂甚至不愿从父那里求得天使大军的援助，且责备了伯多禄的剑。诚然，祂也承认祂的\BibleQuote{心灵忧愁，甚至死亡}，并且肉身软弱；但祂这样做的目的，首先是为使祂自己拥有心灵的忧愁和肉身的软弱，借此向你显示祂内的两种性体都是真正人性的；免得你像如今某些人所引入的那样，认为基督的肉身或灵魂与我们的不同；其次，借着展示它们的状态，使你确信它们离了圣神自己毫无能力。为此，祂先提到\BibleQuote{心灵固然愿意}，好让你注视两种性体各自的本性，从而看出你内心既有心灵的力量，也有肉身的软弱；甚至从中学习该做什么、靠什么去做、以及将什么置于什么之下——即让软弱的服从坚强的，免得你像如今你的作风那样，借口肉身的软弱，反而忽视心灵的力量。祂也曾向父祈求，若是可能，愿那苦难之杯离开祂。你也当如此祈求，但像祂那样，坚持你的立场——只发出恳求，并加上另一句话：\BibleQuote{但不要照我的意愿，而要照你的意愿。} 然而，当你逃跑时，你如何作此祈求？那时你已将挪开苦杯的事握在自己手中，不是做你父所愿意的，而是做你自己所愿意的。\par

9. 宗徒的教导当然在每件事上符合天主的旨意：他们没有忘记也没有省略福音的任何内容。那么，你从哪里证明他们更新了从一城逃往另一城的命令呢？事实上，他们绝不可能规定任何与自己的榜样如此截然相反的事作为逃跑的命令，因为正是从束缚，或从他们因宣认基督徒之名而非因逃避而被囚禁的岛屿，他们写信给教会。\Person{保禄}嘱咐我们扶持软弱者，但最肯定的是，不是在他们逃跑的时候。因为不在场的你们如何扶持他们呢？通过容忍他们？嗯，他说，如果有人在信德软弱上犯了过失，就必须扶持他们，正如（他吩咐）我们应安慰心灰意冷的人；然而，他并没有说他们该被放逐。但当他敦促我们不要给邪恶留地步时，他并未建议我们逃跑，他只教导要克制情欲；若他说要赎回光阴，因为日子邪恶，他愿意我们延长生命，不是通过逃跑，而是通过智慧。此外，他嘱咐我们像光明之子一样照耀，并没有嘱咐我们像黑暗之子一样躲藏。他命令我们坚定不移，当然不是通过逃跑来行相反的事；并且要束腰，不是扮演逃兵或违背福音。他还指出了武器，那些打算逃跑的人是不需要的。在这些武器中，他也提到了盾牌，使你能熄灭魔鬼的箭，当你无疑地抵抗他，并全力承受他的攻击时。因此，\Person{若望}也教导我们必须为弟兄舍命；更何况为主呢。这不能由那些逃跑的人实现。最后，鉴于他自己的默示录，他在其中听到了恐惧者的命运，所以凭个人知识说话，他警告我们必须除去恐惧。他说：\BibleQuote{没有恐惧}，他说，\BibleQuote{但完美的爱除去恐惧，因为恐惧含有刑罚}——无疑是火湖的火。\BibleQuote{那恐惧的人在爱内尚未成全}——即天主的爱。然而，除了那恐惧的人，谁会逃避迫害呢？谁会恐惧呢？除了那不爱的人？是的；你若请教圣神，难道祂比圣神的这句话更赞同什么吗？因为，的确，它几乎激励所有人都去献身殉道，而不是逃避它；所以我们也要提及它。他说，如果你遭受公众羞辱，那对你有益；因为那未在人前受辱的，在主前也必如此。不要羞愧；正义将你带到众人面前。你为何羞于获得荣耀呢？当你在人眼前时，机会就给了你。同样，在别处：不要寻求死在婚床上，或流产中，或温和的热病中，而要殉道而死，为使那为你受苦者得荣耀。\par

10. 但是，有些人毫不在意天主的劝诫，却更愿意将那句世俗智慧的希腊短诗应用于自己身上：\Quote{逃跑的人还会再战；也许还会在战场上再次逃跑。}而且，那个作为逃亡者而被打败的人，何时才能成为胜利者呢？他向他的统帅基督提供了一名配得上的士兵，这士兵受到宗徒如此充分的武装，一听到迫害的号角，就从迫害之日逃跑。我也要引用一句来自世俗的话作为回答：\Quote{死亡真是如此悲惨的事吗？}无论如何，他都必须死，要么作为失败者，要么作为胜利者。但是，尽管他因否认而屈服，他却曾面对并搏斗于酷刑。我宁愿成为被怜悯的对象，也不愿成为令人羞愧的对象。在战场上被标枪刺伤的士兵，比作为逃亡者而皮肤完好的人更光荣。你害怕人吗，基督徒？——你本应被天使所畏惧，因为你要审判天使；你本应被邪魔所畏惧，因为你已经得到了胜过邪魔的能力；你本应被整个世界所畏惧，因为世界也要由你审判。你披戴了基督，你却在魔鬼面前逃跑，因为你已受洗归于基督。在你内的基督被藐视，当你将自己交回给魔鬼，在他面前成为逃亡者。但是，既然你从主那里逃跑，你嘲笑所有逃亡者的徒劳。某位勇敢的先知也曾逃离主，他从\Place{约培}渡海往\Place{塔尔索}方向，仿佛他也能轻易将自己从天主那里转移走；但我发现他，我并非说在海中或陆上，而是在兽腹中，被囚其中三天，既不能寻死，也不能这样逃脱天主。那个人的行为好得多：他尽管害怕天主的仇敌，却不逃避，反而轻看他，依靠主的保护；或者，如果你愿意，他对天主的敬畏更大，因为他站在祂面前，他说：\Quote{祂是主，祂是大能的。万有都属于祂；无论我在哪里，我都在祂手中：愿祂按祂的旨意行事，我不离开；如果祂乐意我死，让祂自己毁灭我，而我为祂保全自己。我宁愿因祂的旨意而死将怨恨加于祂，也不愿因我自己的愤怒而逃脱。}\par

11. 所以，天主的每个仆人都应当这样感受和行动，即使他处于较低的职位，好能藉着忍受迫害而获得晋升，从而得到一个更重要的职位。但是，当身居权位的人——我指的是执事、司铎和主教本人——都选择逃跑时，一个平信徒又如何能明白\BibleQuote{从一城逃到另一城}这句话的真正用意呢？同样，当领袖转身逃跑时，平信徒中又有谁希望能说服人们坚定不移地投入战斗呢？毫无疑问，善牧为羊舍命，正如\Person{梅瑟}的话所预表——那时主基督尚未显现，却已在他身上预表出来：他说：\BibleQuote{如果祢消灭这人民，求祢也把我一同消灭。}但基督亲自证实这些预象，接着说：\BibleQuote{恶牧是那看见狼来就逃跑，撇下羊任凭狼撕碎的人。}这样的牧人将从农场被赶走；他本应于解雇时得到的工资将被扣下作为赔偿；不仅如此，甚至他过去的积蓄也要用来补偿主人的损失，因为\BibleQuote{凡有的，还要给他；凡没有的，连他所有的也要夺去。}因此\BibleBook{}{匝加利亚}警告说：\BibleQuote{刀剑哪，兴起攻击我的牧人，击打羊群；我要伸手攻击牧人。}\BibleBook{}{厄则克耳}和\BibleBook{}{耶肋米亚}也以类似的威胁痛斥他们，因为他们不仅邪恶地吃羊——他们喂养自己而不喂养受托照管的人——而且还驱散羊群，使羊群没有牧人，成为田野所有野兽的猎物。这种情况最常发生在迫害期间教会遭圣职人员遗弃之时。如果有人也认出圣神，他将听见圣神谴责逃跑者。但是如果看守羊群的人在狼入侵时逃跑是不相宜的——不，如果那是绝对不合法的（因为那宣告这类牧人是恶牧的，无疑已经定了他的罪；凡被定罪的，当然就成为不合法）——基于此，被委派管理教会的人在迫害时期就没有逃跑的责任。但另一方面，如果羊群应该逃跑，羊群的监督就没有理由坚守岗位，因为在这种情况下他这样做，无异于毫无理由地给予羊群保护，而羊群由于享有逃跑的自由，并不需要这种保护。\par

\Person{兄弟}，就你所提出的问题而言，我们已经给了你答复和鼓励的意见。但那个询问我们是否应当逃避迫害的人，现在也必须准备好考虑以下问题：如果我们不应当逃避，是否至少可以用钱财赎买自己？因此，超出你的预料，我还要就这一点给你我的建议，明确肯定：我们显然不可逃避的迫害，同样也不该用钱赎买。区别在于支付的方式；但正如逃避是一种不花钱的赎买，赎买也是一种金钱的逃避。确实，这里你也看到了恐惧的劝告。因为你害怕，所以你赎买自己；于是你逃跑。就你的双脚而言，你站住了；就你所付的钱而言，你跑掉了。为什么？在这站立之中，你逃避了迫害，因你买得了脱离迫害；但用金钱赎回一个已被基督用祂的血赎回的人，这如何配得上天主及其行事的方式？祂没有怜惜自己的儿子，为你而交出祂，使祂为我们成了诅咒，因为\BibleQuote{挂在树上的都是被诅咒的}，——\BibleQuote{祂被牵去如同要祭献的羊，又像羔羊在剪毛者面前，祂不开口；祂将背转向鞭打，不，将面颊转向击打者，没有转脸回避唾污，并被列于罪犯之中，被交付于死亡，甚至被钉死在十字架上}。这一切发生，是为将我们从罪中赎回。太阳将我们救赎的日子让给了我们；地狱将它对我们的权柄转移回来，我们的盟约在天堂；永远的门户被举起，好让光荣的君王、大能的主进入，在将人从地上、甚至从地狱赎回后，使他能到达天堂。那么，我们现在该如何看待那与这位光荣者对抗，不，轻视和亵渎祂以如此巨大代价——实在就是祂最宝贵的血——所获得的财货的人？看来，比起降低价值，逃跑更好，如果一个人不肯为自己付出基督为他付出的代价。主确实将他从统治世界的天使权势、从邪恶的鬼神、从这生命的黑暗、从永罚、从永死中赎回。但你却用告密者、士兵或某个卑劣的强盗般的统治者来为他讨价还价——正如他们所说，在衣襟的褶子里——好像他是基督在全世界面前买下并释放的赃物。那么，你愿意用任何价格估价这个自由人，并用任何价格拥有他，但唯独不是用我们所说的主所付出的代价，即祂自己的血吗？（如果不是，）那么你为什么在基督所居住的人身上购买基督，好像祂是某种人类财产？连西满也曾试图这样做，他向宗徒们要钱买基督的圣神。因此，这个人在赎买自己时也买了基督的圣神，他将听到那句话：\BibleQuote{愿你的银子和你一同丧亡！因为你以为天主的恩赐可以用钱买得。}然而，谁又因他是（什么）——一个否认者——而鄙视他呢？因为那勒索者说什么？给我钱：无疑为的是不把他交出来，因为他除了要给你的钱之外，不试图卖给你任何东西。当你把钱放进他手里时，你显然希望\Emph{不}被交出来。但不被交出来，你是否要被公开羞辱？那么，当你不愿意被交出来时，你就不愿意这样被暴露；藉着你的不愿意，你否认了你是不愿意被公开的那一位。不，你说：当我不愿意被公开为我现在这样时，我已经承认了我是不愿意这样被公开的那一位，即一个基督徒。那么，基督能主张你作为祂的见证人，坚定地显扬了祂吗？那赎买自己的人并未这样做。在\Emph{一个人}面前，我确信可以说：你宣认了祂；同样，因你不愿意在许多人面前宣认祂，你就否认了祂。或许，主曾做过你以为你将要做的事？事实上，祂在若瑟的事上确实做了：若瑟不是被他的兄弟们卖的吗？但那是基督的预像，不是为自己出卖的论据；因为若瑟不是被自己出卖，而是被他人，而且是卖给了外邦人，不是作为基督徒，而是作为奴隶。你却要把自己卖给外邦人，好成为基督徒？你改变了预像：你成了恶主人的若瑟。但你能从中得到什么好处？我要再说一遍：一个人的安全本身会宣告，他在逃避迫害时已经跌倒。因此，凡渴望逃脱的人已经跌倒。拒绝殉道就是否认。基督徒靠财富得到保存，并为此拥有财宝，使他不受苦，同时他要在天主面前富有。但事实上，基督为他流了丰富的血。因此，\BibleQuote{贫穷的人是有福的，因为天国是那些只存积灵魂的人的。}如果我们不能服事天主又服事玛门，那么我们能被天主和玛门两者赎回吗？因为谁会比被玛门赎买的人更服事玛门呢？最后，你凭什么例子来证明你可以用金钱避免被出卖？宗徒们处理事情时，在迫害的患难中，何时用金钱解救过自己？他们确实有钱，来自卖地的价钱，放在他们脚前，毫无疑问，在信的人中有许多富人，包括男人和女人，他们也常照顾宗徒的安舒。敖尼息摩、阿桂拉或斯德望何时在受迫害时给过他们这种帮助？保禄确实，当总督斐理斯希望从门徒那里得到钱为他付赎金，他私下也与宗徒商议此事，但保禄既没有自己付，门徒也没有为他付。那些门徒，无论如何，他们哭泣，因为保禄同样固执地决定去耶路撒冷，并忽略一切保护自己免受预言那里将要发生的迫害的手段，最后他们说：\BibleQuote{愿主的旨意成就。}那旨意是什么？无疑他应为主的名受苦，而不是被赎买。因为正如基督为我们舍弃了生命，我们也当为祂舍弃生命；不仅为主自己，而且为我们的弟兄因祂的缘故。这也是若望的教导，他宣告，我们不是要为弟兄付钱，而是要为弟兄舍命。你所不应做的事，无论是赎\Emph{买}一个基督徒，还是\Emph{购买}一个，都没有区别。因此天主的旨意与此一致。看那条件——当然是天主安排的条件，王的心在祂手中——就是列国和帝国。为了增加国库，每天都提供许多手段——财产登记、实物税、自愿捐献、金钱税；但至今从未有过通过使基督徒为个人和教派缴纳赎金来提供任何类似的东西，尽管从如此众多的人数中可以获取巨大的收益，没有人不知道他们。用血买来，用血付清，我们不为我们的头欠任何金钱，因为基督是我们的头。基督不应让我们花钱。如果通过贡金我们为教派的自由付钱，那么殉道又如何能发生以光荣主呢？因此，凡要求以价格获得它的人，就是反对天主的安排。既然凯撒从未以这种方式对我们加征贡赋——事实上，这样的加征永远不可能——而假基督已近在眼前，贪婪基督徒的血，而不是他们的钱——那么，怎能向我指出那条命令：\BibleQuote{凯撒的归凯撒？}一个士兵，无论是告密者还是敌人，以威胁从我这里勒索金钱，并没有代表凯撒索取任何东西；相反，当他因贿赂而放我走时——我作为基督徒，按人的法律是罪犯——他在做相反的事。我欠凯撒的是另一种\Emph{\OriginalTerm{德纳留}{denarius}}，它是属于他的东西，当时的问题正是关于它，那是一种贡金硬币，确实应由该纳税的人缴纳，而不是由儿女。或者，我该如何将天主的归给天主？——当然，因此那是祂自己的肖像和刻有祂名字的金钱，也就是一个基督徒的人？但我欠天主什么，如同我欠凯撒\Emph{\OriginalTerm{德纳留}{denarius}}一样，岂不就是祂自己的儿子为我所流的血吗？现在，如果我确实欠天主一个人，即我自己的血，但我正处于这样的关头：债务被要求偿还，那么我若不尽力偿还，无疑是犯了欺哄天主的罪。我很好地遵守了诫命：如果我将凯撒的归还给凯撒，却拒绝将天主的归还给天主！\par

13. 然而，凡向我求乞的，我也将因爱德施予，决非出于任何胁迫。\Quote{谁求？}祂说。但施胁迫者并非在求。以威胁强索而不得者，并非渴求，而是强迫。他前来并非为求怜悯，而是令人畏惧，故他所寻的并非施舍。因此，我将施予，是因怜悯而非因恐惧，当受者尊崇天主并为我祝福时；而非当他反以为对我施惠，又目睹其掠夺物而说：\Quote{罪钱}时。难道我甚至要恼恨仇敌吗？但仇恨也有其他缘由。然而祂并没有说\Quote{背叛者}或\Quote{迫害者}或\Quote{以威胁恐吓你的人}。因为如果不以金钱赎身，我将更把这炭火堆在此人头上！耶稣说：\Quote{同样，}祂说，\Quote{拿走你外衣的，连内衣也让他拿去。}但那是指那试图夺我财物的人，而非我的信德。内衣我也愿给，只要不受背叛的威胁。若他威胁，我甚至要索回外衣。就连如今，主的声明也有其自身的理由与规律，并非无限或普遍适用。因此祂吩咐我们给每一个求者，而祂自己却不给求征兆的。否则，若你主张应当不加区别地给一切求者，那么在我看来，你似乎也要给发烧者酒，甚至给寻死者毒药或刀剑。但如何理解\Quote{要用不义的钱财结交朋友}？让前面的比喻教导你吧。这话是对犹太人说的；因为他们受托管理主的家业不善，本应从玛门之人（即我们）中为自己预备朋友而非仇敌，并将我们从阻隔天主的罪债中解救出来——若他们将祝福施予我们，按主所给的理由，即当恩宠开始离开他们时，他们可转向我们的信仰，被接入永恒的居所。现在，你若喜欢对这比喻和警句作别的解释，只要清楚看出一点：我们的反对者若我们以玛门结交他们，他们断不会接我们进入永恒住所。但怯懦不能使人相信什么呢？仿佛圣经既允许他们逃跑，又命令他们赎身！最终，若一人或另一人以此得救，仍不足够。整个教会已为自己\Quote{全体地}课税。我不知道这是可悲还是可耻：在商贩、扒手、浴场窃贼、赌徒和皮条客中，基督徒竟也被列为纳税者，编入自由士兵和探子的名册。难道宗徒们如此有远见地设立了这类监督的职务，使其在位者能借为羊群提供（同样自由）的借口，而无忧无虑地享受其管辖？诚然，基督在回归父前，吩咐用类似农神节中赠礼的方式从士兵那里购买这种和平！\par

14. 但你却说：\Quote{我们如何聚集呢？如何遵守主的规诫呢？}当然，正如宗徒们那样，他们以信德而非金钱护卫；这信德若能移山，更能移走一个士兵。让你的保障是智慧，而非贿赂。因为即使你买通了士兵的干预，也未必能立即从民众那里得到完全的安全。因此，你需要的一切保障就是拥有信德和智慧：你若不用这些，甚至会失去自己赎来的解脱；你若运用它们，就无需任何赎金。最后，若白天不能聚集，你有黑夜，基督的光明照耀其黑暗。你不能逐个地奔走在他们中间。以三人教会为足。有时你宁可不看见你的群众，也不该使自己（受贡赋束缚）。为基督保持祂许配的童贞女纯洁，无人可从中牟利。兄弟啊，这些事于你或许显得严苛而不可忍受，但请记住天主曾说过：\Quote{能领受的，就让他领受}，意思是，不能领受的，就让他走自己的路。害怕受苦的人，不属于那受苦的主。但不怕受苦的人，将在爱德中成全——即对天主的爱德；\Quote{因为完美的爱德驱除恐惧}。\Quote{为此，许多蒙召，但少数被选。}问题不在于谁愿走宽路，而在于谁走窄路。因此，那护慰者是必需的，祂引导人进入一切真理，并激励人坚忍一切。那些接受了祂的人，既不会屈身逃跑，也不会花钱买脱，因为他们有主自己，祂将与我们同在，在我们受苦时帮助我们，并在我们受审时作我们的口。\par

\SourceRecord{Translated by S. Thelwall.；Ante-Nicene Fathers；Vol. 4.；Edited by Alexander Roberts, James Donaldson, and A. Cleveland Coxe.；Buffalo, NY: Christian Literature Publishing Co.,；1885.；<http://www.newadvent.org/fathers/0409.htm>.}
